\documentclass{article}

% if you need to pass options to natbib, use, e.g.:
    \PassOptionsToPackage{numbers, compress}{natbib}
% before loading neurips_2026

% The authors should use one of these tracks.
% Before accepting by the NeurIPS conference, select one of the options below.
% 0. "default" for submission
\usepackage{neurips_2026}
% the "default" option is equal to the "main" option, which is used for the Main Track with double-blind reviewing.
% 1. "main" option is used for the Main Track
%  \usepackage[main]{neurips_2026}
% 2. "position" option is used for the Position Paper Track
%  \usepackage[position]{neurips_2026}
% 3. "eandd" option is used for the Evaluations & Datasets Track
 % \usepackage[eandd]{neurips_2026}
 % if you want to opt-in for a de-anomymized submission (not recommended, but OK):
 %\usepackage[eandd, nonanonymous]{neurips_2026}
% 4. "creativeai" option is used for the Creative AI Track
%  \usepackage[creativeai]{neurips_2026}
% 5. "sglblindworkshop" option is used for the Workshop with single-blind reviewing
 % \usepackage[sglblindworkshop]{neurips_2026}
% 6. "dblblindworkshop" option is used for the Workshop with double-blind reviewing
%  \usepackage[dblblindworkshop]{neurips_2026}

% After being accepted, the authors should add "final" behind the track to compile a camera-ready version.
% 1. Main Track
 % \usepackage[main, final]{neurips_2026}
% 2. Position Paper Track
%  \usepackage[position, final]{neurips_2026}
% 3. Evaluations & Datasets Track
 % \usepackage[eandd, final]{neurips_2026}
% 4. Creative AI Track
%  \usepackage[creativeai, final]{neurips_2026}
% 5. Workshop with single-blind reviewing
%  \usepackage[sglblindworkshop, final]{neurips_2026}
% 6. Workshop with double-blind reviewing
%  \usepackage[dblblindworkshop, final]{neurips_2026}
% Note. For the workshop paper template, both \title{} and \workshoptitle{} are required, with the former indicating the paper title shown in the title and the latter indicating the workshop title displayed in the footnote.
% For workshops (5., 6.), the authors should add the name of the workshop, "\workshoptitle" command is used to set the workshop title.
% \workshoptitle{WORKSHOP TITLE}

% "preprint" option is used for arXiv or other preprint submissions
 % \usepackage[preprint]{neurips_2026}

% to avoid loading the natbib package, add option nonatbib:
%    \usepackage[nonatbib]{neurips_2026}

\usepackage[utf8]{inputenc} % allow utf-8 input
\usepackage[T1]{fontenc}    % use 8-bit T1 fonts
\usepackage{hyperref}       % hyperlinks
\usepackage{url}            % simple URL typesetting
\usepackage{booktabs}       % professional-quality tables
\usepackage{amsfonts}       % blackboard math symbols
\usepackage{nicefrac}       % compact symbols for 1/2, etc.
\usepackage{microtype}      % microtypography
\usepackage{xcolor}         % colors
\usepackage{graphicx}
\usepackage{subcaption}
\usepackage{wrapfig}
\usepackage{xspace}
% \newcommand{\pgmstruct}{\textsc{PGM-Struct}\xspace}
% \newcommand{\pgmdecision}{\textsc{PGM-Decision}\xspace}
\newcommand{\bench}{\textsc{ACED-Bench} }
\newcommand{\pgmstruct}{\textsc{ACED-Struct} }
\newcommand{\pgmdecision}{\textsc{ACED-Decision} }

% Note. For the workshop paper template, both \title{} and \workshoptitle{} are required, with the former indicating the paper title shown in the title and the latter indicating the workshop title displayed in the footnote. 
% \title{A Neuro-Symbolic Benchmark for AI Co-Scientists: \\
% Active Causal Discovery with LLM Agents in Probabilistic Graphical Environments}
\title{\textsc{ACED-Bench}: Evaluating How LLM Agents Acquire and Act on Causal Evidence}


% The \author macro works with any number of authors. There are two commands
% used to separate the names and addresses of multiple authors: \And and \AND.
%
% Using \And between authors leaves it to LaTeX to determine where to break the
% lines. Using \AND forces a line break at that point. So, if LaTeX puts 3 of 4
% authors names on the first line, and the last on the second line, try using
% \AND instead of \And before the third author name.


\author{%
  David S.~Hippocampus\thanks{Use footnote for providing further information
    about author (webpage, alternative address)---\emph{not} for acknowledging
    funding agencies.} \\
  Department of Computer Science\\
  Cranberry-Lemon University\\
  Pittsburgh, PA 15213 \\
  \texttt{hippo@cs.cranberry-lemon.edu} \\
  % examples of more authors
  % \And
  % Coauthor \\
  % Affiliation \\
  % Address \\
  % \texttt{email} \\
  % \AND
  % Coauthor \\
  % Affiliation \\
  % Address \\
  % \texttt{email} \\
  % \And
  % Coauthor \\
  % Affiliation \\
  % Address \\
  % \texttt{email} \\
  % \And
  % Coauthor \\
  % Affiliation \\
  % Address \\
  % \texttt{email} \\
}


\begin{document}


\maketitle


\begin{abstract}
Large language models are increasingly used as scientific assistants, but it remains difficult to evaluate how agents acquire, analyze, and commit to causal evidence under controlled interventional ground truth.
We introduce \bench, a neuro-symbolic benchmark for evaluating LLM agents as sample-based causal investigators.
Each task is grounded in a hidden probabilistic graphical model and presented through a natural-language scientific scenario.
Agents may request observational or interventional samples, write code to analyze finite data, and decide what evidence to collect next under a fixed budget.

\bench contains \pgmstruct, 180 structural questions testing statistical and causal queries, and \pgmdecision, 60 causal-decision questions testing intervention choice under safety, robustness, mediation, satisficing, and validity constraints. Worlds are generated by combining LLM-written semantic surface forms with programmatic graph construction, CPD sampling, validation, and oracle answer computation, enabling scalable instance generation without manual labeling. We evaluate seven contemporary LLMs, with two earlier-generation models as references, across four agent designs. Active evidence collection improves structural reasoning across models, while causal-decision performance is more variable, showing that causal decision-making requires choosing the right contrast, not merely accessing more data.

Beyond final accuracy, our evidence-ledger analysis separates failures of evidence acquisition, evidence retention, and stopping after sufficient evidence has been obtained. These results show that reliable AI co-scientists require not only data access and tool use, but also mechanisms for evidence tracking, commitment, and stopping.
\end{abstract}

\section{Introduction}
\label{sec:intro}

Large language models (LLMs) are increasingly used as scientific assistants: they read papers, write code, analyze data, and propose hypotheses.
But if LLMs are to become reliable AI co-scientists, we need to evaluate not only what they answer, but how they investigate.
Scientific investigation requires choosing what evidence would be informative, collecting data, analyzing uncertain observations, revising beliefs, and deciding when the evidence is sufficient.
This is especially important for causal questions, where observational associations can disagree with interventional effects.

Existing benchmarks cover three capabilities needed by AI co-scientists, but rarely combine them: formal causal reasoning, statistical reasoning from data, and interactive scientific tool use.
Causal-reasoning benchmarks evaluate whether models understand concepts such as intervention, confounding, and counterfactuals~\citep{jinCLadderAssessingCausal2023,wangCausalBenchComprehensiveBenchmark2024,jinCanLargeLanguage2024}.
Statistical-reasoning benchmarks evaluate probability estimation, uncertainty, and data analysis~\citep{liuAreLLMsCapable2024,paruchuriWhatAreOdds2024,nafarReasoningUncertainText2025,zhuAreLargeLanguage2024}.
Scientific-agent and data-analysis benchmarks evaluate interactive tool use, data workflows, and multi-step research behavior~\citep{wangScienceWorldYourAgent2022a,jansenDiscoveryWorldVirtualEnvironment2024,majumderDiscoveryBenchDataDrivenDiscovery2024,guBLADEBenchmarkingLanguage2024,hu2024infiagentdabench,hu2024infiagentdabench,laiKramaBenchBenchmarkAI2026}.
Recent interactive causal benchmarks and LLM-guided intervention methods study graph discovery or intervention targeting~\citep{chenAutoBenchAutomatedBenchmark2025,liCanLargeLanguage2025}.
\bench targets a different intersection: agents must choose observational or interventional evidence, estimate statistical and causal quantities from samples, make constrained causal decisions, and preserve the evidence needed to justify the final answer.

We introduce \bench, a neuro-symbolic benchmark for evaluating LLM agents as sample-based causal investigators.
Each task is grounded in a hidden probabilistic graphical model (PGM) with categorical variables, a causal graph, and conditional probability tables.
The evaluator knows the full PGM, while the agent sees only a natural-language scientific scenario, variable metadata, non-intervenability annotations, and a question.
The agent may request observational samples, request interventional samples, write code to analyze finite data, and decide what to test next under a fixed budget.
Because the underlying PGM is known, every accepted world has oracle-checkable answers.
Figure~\ref{fig:benchmark_overview} illustrates the interaction loop for a causal-decision task.

\bench contains two datasets: \pgmstruct, with 180 structural questions testing marginal, conditional, independence, and basic causal reasoning, and \pgmdecision, with 60 causal-decision questions testing intervention choice under safety, robustness, mediation, satisficing, and validity constraints.
Together, they test whether agents can recover distributional and causal facts from actively requested evidence, and whether they can use those facts to choose constrained interventions.

We evaluate seven contemporary LLMs, with two earlier-generation reference models, across four agent designs.
Active evidence collection improves structural reasoning across models, while \pgmdecision remains more sensitive to selecting the right constrained causal contrast.
We further analyze modular code-using trajectories with an evidence ledger that separates evidence acquisition, evidence retention, ignored evidence, and stopping behavior.

Our contributions are:
\begin{itemize}
    \item We introduce \bench, an interactive PGM-based benchmark for evaluating LLM agents on active causal evidence acquisition under exact observational and interventional ground truth, with two datasets: \pgmstruct containing 180 structural statistical and causal questions, and \pgmdecision containing 60 causal-decision questions involving safety, robustness, mediation, satisficing, and invalid-premise constraints.
    
    \item We provide a scalable neuro-symbolic generation pipeline in which LLMs create semantic surface forms while code enforces graph construction, CPD sampling, intervention semantics, validation, and oracle answer computation, enabling scalable instance generation without manual labeling.
    
    \item We evaluate seven contemporary LLMs, with two earlier-generation models as references, across four agent designs on 240 total questions, showing that active evidence collection improves structural reasoning across models while causal-decision tasks remain more sensitive to selecting the right causal contrast.
    
    \item We introduce an evidence-ledger analysis that decomposes agent trajectories into whether agents acquire sufficient evidence, retain it in the final answer, ignore earlier correct evidence, or continue querying after sufficient evidence has been obtained.
\end{itemize}


\begin{figure}[t]
    \centering
    \includegraphics[width=0.5\linewidth]{diagram-advanced.pdf}
    \caption{
    Overview of \bench on a causal-decision task.
    Each task is generated from a hidden probabilistic graphical model (PGM).
    The benchmark oracle has access to the causal graph, conditional probability tables, and exact interventional answers.
    The LLM scientist agent only observes a natural-language scenario, variable metadata, non-intervenability annotations, and the question.
    Under a fixed budget, the agent may request observational or interventional samples, analyze the resulting data, and recommend an intervention that improves the target while satisfying the task constraint.
    }
    \label{fig:benchmark_overview}
\end{figure}

\section{Related Work}
\label{sec:related}

\paragraph{Agentic scientific and data-analysis benchmarks.}
Recent benchmarks increasingly evaluate LLMs as agents that interact with tools, environments, or data sources.
ScienceWorld and DiscoveryWorld test scientific exploration in interactive environments requiring hypothesis formation, experimentation, and explanation~\citep{wangScienceWorldYourAgent2022a,jansenDiscoveryWorldVirtualEnvironment2024}.
DiscoveryBench, BLADE, DA-Bench, and KramaBench evaluate data-driven investigation, code execution, and data-to-insight workflows~\citep{majumderDiscoveryBenchDataDrivenDiscovery2024,guBLADEBenchmarkingLanguage2024,hu2024infiagentdabench,laiKramaBenchBenchmarkAI2026}.
These benchmarks establish that scientific-agent evaluation should include planning, tool use, and iterative interaction.
\bench{} asks a complementary question: whether an agent can gather and use evidence in a controlled causal system.
Agents still interact, query data, and write code, but every world has exact observational and interventional semantics, allowing oracle-checkable evaluation of causal evidence, causal decisions, and the trajectory that produced them.

\paragraph{Causal and statistical reasoning benchmarks.}
Causal-reasoning benchmarks test whether LLMs understand concepts such as association, intervention, confounding, and counterfactuals~\citep{jinCLadderAssessingCausal2023,wangCausalBenchComprehensiveBenchmark2024,jinCanLargeLanguage2024,duIceCreamDoesnt2026}.
Statistical-reasoning benchmarks evaluate probability, uncertainty, calibration, and data analysis~\citep{liuAreLLMsCapable2024,paruchuriWhatAreOdds2024,nafarReasoningUncertainText2025,zhuAreLargeLanguage2024,rendaOpenEstimateEvaluatingLLMs2026,luStatEvalComprehensiveBenchmark2025}.
These benchmarks provide controlled tests of important reasoning skills, but the relevant evidence is usually given directly as text, a graph, or a fixed dataset.
In \bench, evidence is endogenous: the agent must decide which observational or interventional samples to request, analyze finite data, and determine when the evidence is sufficient.

\paragraph{LLMs for causal discovery and experimental design.}
Several works use LLMs to assist causal discovery, for example by supplying background knowledge, orienting edges, reducing equivalence classes, or guiding intervention selection~\citep{longCausalDiscoveryLanguage2023,takayamaIntegratingLargeLanguage2025,abdulaalCausalModellingAgents2024,liCanLargeLanguage2025}.
This line of work is close to our causal motivation, but often treats the LLM as one component in a causal-discovery algorithm.
Interactive causal benchmarks such as Auto-Bench also evaluate iterative discovery, but focus primarily on hidden-graph or adjacency-matrix reconstruction~\citep{chenAutoBenchAutomatedBenchmark2025}.
Our benchmark instead evaluates sample-based causal investigation: agents answer statistical and causal queries, choose constrained interventions, and are scored by exact PGM ground truth.

\paragraph{Active causal learning and intervention design.}
Our setting is grounded in causal inference and probabilistic graphical models~\citep{pearl2009causality,koller2009probabilistic}.
It is also related to active causal learning, where interventions are selected to identify structure or reduce uncertainty~\citep{he2008active,hauserCharacterizationGreedyLearning2012,tothActiveBayesianCausal2022,zhangActiveLearningOptimal2023}.
Classical methods typically assume an explicit statistical model, posterior, or acquisition objective.
In contrast, \bench evaluates general-purpose LLM agents that interact through natural language, sampling APIs, and optional code execution.
The goal is not to propose a new intervention-selection algorithm, but to test whether agents can choose informative evidence and use it correctly.

\paragraph{Trajectory-level evaluation and evidence management.}
Final-answer accuracy can hide important differences between agents that make the same mistake for different reasons.
Process-level evaluation and language-agent trajectory analysis therefore study intermediate reasoning, actions, observations, and tool use~\citep{lightmanLetsVerifyStep2023,yaoReActSynergizingReasoning2023b,shinnReflexionLanguageAgents2023,kimFinalAnswerEvaluating2025,fanAgentProcessBenchDiagnosingStepLevel2026}.
Our evidence ledger specializes this idea to causal investigation.
Each turn is judged by whether it computes the relevant estimand, whether the evidence supports the gold answer, and whether the evidence is retained in the final response.
This lets us separate failures of evidence acquisition, evidence retention, and stopping.

\section{\bench}
\label{sec:benchmark}

\bench evaluates LLM agents as sample-based causal investigators in hidden probabilistic graphical model (PGM) worlds.
Each world is a tuple
\[
    W = (G, \mathcal{V}, \mathcal{P}, s, \mathcal{Q}, \mathcal{N}),
\]
where $G=(V,E)$ is a DAG over categorical variables, $\mathcal{V}$ contains variable names, descriptions, and value sets, $\mathcal{P}$ defines the Bayesian-network CPDs, $s$ is a natural-language scientific scenario, $\mathcal{Q}$ is a set of questions, and $\mathcal{N}\subseteq V$ is the set of non-intervenable variables.
The evaluator has access to the full PGM.
The agent observes only the scenario, variable metadata, non-intervenability annotations, the question, and its interaction history.
At each turn, the agent may request observational samples from $P_W(V)$, request interventional samples from $P_W(V\mid \mathrm{do}(X=x))$, run code when permitted, or submit a final answer.
The agent is evaluated under a fixed budget of successful data queries.
Figure~\ref{fig:benchmark_overview} illustrates this interface.

\paragraph{Neuro-symbolic construction.}
\bench separates the natural-language surface form of a world from the causal system used for scoring.
LLMs write variable names, value descriptions, scientific scenarios, and task text.
Code constructs the graph, samples CPDs, defines intervention semantics, validates the generated world, and computes oracle answers.
Thus, the LLM makes each task readable, but the causal truth is produced and checked programmatically.

This construction makes the benchmark scalable by design.
Additional worlds, graph sizes, domains, and archetype instances can be sampled and validated without manual answer labeling.
The main scaling costs are generation time, validation rejection rate, simulator runtime, and LLM evaluation cost, rather than human construction of gold answers.
Appendix~\ref{app:validation} gives the validation criteria used to accept generated worlds.

\paragraph{\pgmstruct.}
\pgmstruct contains 30 worlds and 180 structural questions.
For each world, we sample a random DAG with controlled graph size, fan-in, and path length, then use a backend LLM to assign a scientific scenario and variable metadata across eight topical domains.
Root CPDs are sampled from Dirichlet priors, and non-root CPDs are parameterized as logistic conditional probability tables.
The resulting questions test marginal probabilities, conditional probabilities, independences, basic causal effects, and set-valued cause/effect relationships.

Before accepting a structural world, we validate the exact distribution induced by its graph and CPDs.
Worlds are rejected when CPD sampling creates accidental independences, unstable dependencies, or ambiguous structural questions.
This ensures that structural questions have oracle-checkable answers rather than artifacts of weak effect sizes.
Appendix~\ref{app:validation} gives the exact structural validation tests.

\paragraph{\pgmdecision.}
\pgmdecision contains 60 worlds and 60 causal-decision questions.
Each world is built around one of five decision archetypes.
\textbf{Safety} asks the agent to improve a target without causing unacceptable harm to a side outcome.
\textbf{Satisficing} asks the agent to find any intervention that improves the target by a meaningful amount, rather than the single best intervention.
\textbf{Subgroup robustness} asks for an intervention that improves outcomes across all subgroup values, not only on average.
\textbf{Mediation} asks whether an effect operates through a proposed mediator or which candidate mediator explains the effect.
\textbf{Invalid premise} asks the agent to reject a proposed intervention that is not causally valid, such as intervening on a non-intervenable variable, and to recommend a valid alternative when one exists.
These archetypes require agents to identify the relevant causal contrast before analyzing data.

Figure~\ref{fig:decision_generation} summarizes the \pgmdecision generation pipeline.
The decision pattern controls the semantic roles, graph constraints, validation checks, and oracle answer.
The LLM supplies only the surface form.
This lets \pgmdecision contain natural scenarios with non-obvious causal behavior, while preventing the LLM from determining the ground truth.
Each world is accepted only if exact interventional inference confirms the intended archetype signature; worlds are rejected when graph constraints fail, effect margins are too small, or the oracle answer is ambiguous.
Appendix~\ref{app:validation} gives the archetype-specific validation rules and thresholds.

\begin{figure}[t]
    \centering
    \includegraphics[width=0.5\linewidth]{PGM-Decision.pdf}
    \caption{
    Generation pipeline for \pgmdecision.
    A decision pattern and topic determine semantic roles, graph constraints, CPD construction, validation checks, and the final question.
    LLMs provide readable scientific surface forms, while code enforces causal validity and computes oracle answers.
    }
    \label{fig:decision_generation}
\end{figure}

\paragraph{Dataset summary and gold answers.}
Table~\ref{tab:dataset_summary} summarizes the two datasets.
Both span eight topical domains: Criminal Justice, Education, Hospital Data, Labor \& Policy, Screening \& Diagnosis, Social Science, Treatment Effectiveness, and User Behavior.

\begin{table}[t]
\centering
\small
\caption{
Dataset summary.
\pgmstruct evaluates structural statistical and causal reasoning; \pgmdecision evaluates causal decisions under scientific constraints.
}
\label{tab:dataset_summary}
\begin{tabular}{lcccc}
\toprule
Dataset & Worlds & Questions & Graph sizes & Questions per world \\
\midrule
\pgmstruct   & 30 & 180 & $n\in\{10,20,30\}$ & 6 \\
\pgmdecision & 60 & 60  & $n\in\{10,15\}$    & 1 \\
\bottomrule
\end{tabular}
\end{table}

All gold answers are computed from the underlying PGM.
Structural answers use graph reachability, d-separation, and exact probabilistic queries.
Decision answers are computed by enumerating candidate interventions and evaluating their exact interventional distributions.
Generated worlds are filtered using effect-size margins, threshold gaps, and tie tolerances to avoid unstable labels.
Appendix~\ref{app:structural_category_breakdown} reports the fine-grained structural question breakdown.

\section{Agent and Evaluation Protocol}
\label{sec:protocol}

All agents receive the same scenario, variable metadata, non-intervenability annotations, question, and interaction history.
They differ only in how they choose evidence and whether they can execute code.
A fixed world-parser LLM maps natural-language data requests to typed simulator queries.
The parser is not evaluated as a scientist agent: it performs no causal reasoning, and all query validity checks are deterministic.
Holding this interface fixed ensures that model comparisons reflect evidence selection and analysis rather than differences in data access.

\paragraph{Agent variants.}
We evaluate four agent designs.
The \textbf{guess-shot} agent answers immediately without data access.
The \textbf{reasoning} agent requests observational or interventional samples and reasons from simulator summaries, but cannot execute code.
The \textbf{coder} agent can additionally write Python code to analyze sampled CSV files.
The \textbf{modular coder} agent separates planning, code analysis, interpretation, and experiment design into explicit stages.
These variants are not intended to exhaust the design space; they isolate the effects of active sampling, code execution, and modular decomposition under the same benchmark interface.

\paragraph{Models.}
We evaluate seven contemporary scientist models: Claude Opus-4.6, Claude Sonnet-4.6, Gemini-3-Pro, DeepSeek-V4-Pro, Qwen3.6-27B, GPT-OSS-120B, and Llama-4-Maverick.
We also include GPT-4o and Llama-3-70B as earlier-generation reference models.
For each model, we run the agent variants supported by our inference stack on both datasets.
Available modular-coder trajectories are additionally used for the evidence-ledger analysis in Section~\ref{subsec:evidence_ledger_results}.

\paragraph{Scoring.}
Binary questions are scored by accuracy.
Set-valued cause/effect questions are scored by exact match and $F_1$.
Categorical mediator questions are scored by exact label match.
Action-choice questions are parsed into a variable-value intervention and scored against the oracle effect table.
We report tie-accepted accuracy unless otherwise stated.
For satisficing tasks, the predicted intervention must also satisfy the stored improvement threshold.

\paragraph{Uncertainty estimates.}
For \pgmstruct, we report world-clustered bootstrap intervals, resampling worlds rather than individual questions to preserve the dependence among questions generated from the same graph.
For \pgmdecision, each world contributes one question, so world-level and question-level bootstraps coincide.
For within-world comparisons between agent variants, we use paired bootstrap intervals.

\paragraph{Behavioral metrics.}
In addition to final-answer accuracy, we log query budget use, failed queries, code executions, code failures, give-up rate, and binary label bias.
These metrics help distinguish failures of causal reasoning from failures of tool use, parsing, or budget management.

\subsection{Evidence-Ledger Analysis}
\label{subsec:ledger_method}

Final-answer accuracy does not reveal why an interactive agent failed.
A trajectory may fail because the agent never collected the right evidence, because it collected the evidence but did not use it in the final answer, or because it kept querying after sufficient evidence had already been obtained.
We therefore introduce an evidence ledger for diagnostic trajectory evaluation.

For each turn in a modular code-using trajectory, an LLM annotator judges whether the turn computed the correct estimand for the question, whether the resulting evidence supports the gold answer, and whether the evidence is adequate.
A turn is marked as \emph{strong evidence} if all three conditions hold.
From these turn-level labels, we compute four trajectory-level diagnostics:
\textbf{evidence acquisition}, whether the trajectory ever obtains strong evidence;
\textbf{evidence retention}, whether the final answer is correct conditional on obtaining strong evidence;
\textbf{ignored evidence}, whether a failed trajectory nevertheless contained strong evidence; and
\textbf{extra turns after evidence}, the number of turns after the first strong-evidence turn.

Because ledger labels depend on trajectory annotation, we treat the ledger as diagnostic rather than as the primary benchmark score.
The annotator receives the question, gold answer, interaction history, code outputs, and final response, but its judgments may still contain errors.
We therefore use the ledger to interpret failure mechanisms, while final-answer accuracy remains the main quantitative metric.
Full metric definitions are provided in Appendix~\ref{app:ledger_definitions}.

\section{Results}
\label{sec:results}

\subsection{Active evidence collection drives most of the gain}
\label{subsec:overall_results}

Figure~\ref{fig:overall_accuracy} summarizes final-answer accuracy across datasets, models, and agent designs.
On \pgmstruct, all evaluated models improve substantially when allowed to query the simulator rather than answering from the scenario alone.
Across models, active sampling improves substantially over guess-only answering.
This shows that the benchmark is not primarily measuring story-level priors: agents benefit from gathering evidence from the hidden PGM.

The strongest structural results come from code-using or modular agents.
Gemini-3-Pro reaches $96.1\%$ with active code, Opus-4.6 reaches $94.8\%$ with the modular agent, and Qwen3.6-27B reaches $91.7\%$ with the modular agent.
However, code is not uniformly beneficial.
For GPT-4o, Llama-3-70B, Llama-4-Maverick, and GPT-OSS-120B, the single-prompt coder agent can underperform active sampling.
The most extreme case is GPT-OSS-120B on \pgmstruct, where active sampling reaches $76.7\%$, the coder agent drops to $43.3\%$, and the modular coder recovers to $80.0\%$.
This suggests that code helps only when the agent can maintain the experiment loop reliably; otherwise, decomposing planning, analysis, interpretation, and next-experiment selection can be more important than code access itself.

Causal-decision tasks produce larger separation among models.
On \pgmdecision, the best code-using cells are near ceiling for several current models: Opus-4.6 reaches $100\%$ with active code, Sonnet-4.6 reaches $97.9\%$ with the modular agent, DeepSeek-V4-Pro reaches $95.0\%$ with active code, and Gemini-3-Pro and Qwen3.6-27B each reach $91.7\%$ with active code.
By contrast, GPT-OSS-120B, GPT-4o, and Llama-3-70B remain substantially lower.
Thus, \pgmstruct tests whether agents can recover hidden statistical and causal relationships from samples, while \pgmdecision more sharply tests whether they can choose and act on the right constrained causal comparison.

\begin{figure}[t]
    \centering
    \includegraphics[width=0.5\linewidth]{fig1_combined_accuracy.pdf}
    \caption{
    Overall accuracy on \pgmstruct and \pgmdecision.
    Active sampling improves over guess-only answering for every evaluated model.
    Code and modular decomposition help the strongest models, but single-prompt code can fail for models that struggle to maintain the long agent loop.
    Error bars are $95\%$ world-clustered bootstrap intervals.
    }
    \label{fig:overall_accuracy}
\end{figure}

\subsection{Causal-decision tasks stress constrained causal comparison}
\label{subsec:advanced_results}

Figure~\ref{fig:advanced_archetypes} breaks down \pgmdecision accuracy by archetype.
The easiest archetypes are not always the most causally diagnostic.
Subgroup-robust and satisficing tasks are often near saturation for the strongest models, and subgroup-robust questions are also relatively answerable from the story alone.
By contrast, safety-constrained and mediator-structure tasks preserve more spread across the model fleet.

The safety-constrained archetype is especially diagnostic because a correct answer must improve the target while also verifying that the same intervention does not harm a side outcome.
Several current models perform well on this archetype, but weaker/reference models drop sharply.
GPT-4o reaches only $13.3\%$ on safety-constrained questions in the plotted code-based evaluation, and Llama-3-70B also remains weak.
Mediator-structure questions are similarly discriminative because the agent must identify the relevant mediated or direct causal contrast, rather than merely estimate a single treatment effect.

Because each archetype contains only about $10$--$15$ questions, individual cells should be interpreted as diagnostic rather than definitive.
The robust pattern is that \pgmdecision differentiates models most clearly when the task requires selecting the right causal contrast under an additional scientific constraint.

\begin{figure}[t]
    \centering
    \includegraphics[width=0.8\linewidth]{fig2_advanced_archetypes.pdf}
    \caption{
    Accuracy by causal-decision archetype.
    The plotted cells use the current code-based runs for the newer models and the available code-based reference runs for GPT-4o and Llama-3-70B.
    Safety-constrained and mediator-structure tasks preserve the most diagnostic spread, suggesting failures in formulating the constrained causal comparison rather than merely estimating an already-specified effect.
    }
    \label{fig:advanced_archetypes}
\end{figure}

\subsection{The evidence ledger separates acquisition, retention, and stopping}
\label{subsec:evidence_ledger_results}

Final-answer accuracy does not reveal why an interactive agent failed.
A trajectory can fail because the agent never obtains the right evidence, because it obtains strong evidence but does not preserve it in the final answer, or because it continues querying after sufficient evidence has already been found.
We therefore analyze modular code-using trajectories with the evidence ledger defined in Section~\ref{subsec:ledger_method}.
Figures~\ref{fig:ledger_scatter} and~\ref{fig:ledger_bars} summarize the resulting diagnostics.

\paragraph{Evidence acquisition is the main axis of variation.}
Across both datasets, evidence retention is usually high once strong evidence appears.
On \pgmstruct, all models except Llama-3-70B retain strong evidence at rates of at least $85\%$, and the strongest models are near or above $97\%$.
On \pgmdecision, retention also remains high for the modular trajectories we annotate.
The larger variation is in evidence acquisition: whether the agent ever chooses the right estimand and obtains evidence that supports the gold answer.
On \pgmstruct, acquisition ranges from $57.5\%$ for GPT-4o to $87.9\%$ for Opus-4.6.
On \pgmdecision, the spread is sharper, ranging from $33.3\%$ for Llama-4-Maverick to $100\%$ for Opus-4.6.
This supports the main process-level conclusion: for current agents, the hardest part is often not using correct evidence once it appears, but deciding what evidence to obtain in the first place.

\begin{wrapfigure}{r}{0.4\linewidth}
    \vspace{-0.8em}
    \centering
    \includegraphics[width=\linewidth]{fig3_ledger_scatter_clean.pdf}
    \vspace{-0.8em}
    \caption{
    Evidence acquisition versus evidence retention for modular code-using trajectories.
    Filled markers show \pgmstruct; open markers show \pgmdecision.
    Bubble size indicates ignored evidence.
    }
    \label{fig:ledger_scatter}
    \vspace{-1em}
\end{wrapfigure}

\paragraph{Ignored evidence reveals distinct failure modes.}
Some failures occur even after the trajectory has already contained strong evidence.
These cases are captured by the ignored-evidence diagnostic.
For example, Opus-4.6 on \pgmdecision has perfect evidence acquisition, but its few failures occur despite earlier strong evidence.
Sonnet-4.6 shows the opposite pattern on \pgmdecision: its ignored-evidence rate is zero, meaning that its rare failures are acquisition failures rather than retention failures.
These failure denominators are small, so the exact percentages should not be over-interpreted.
The important point is qualitative: two models with similar final accuracy can fail for different reasons, suggesting different interventions such as better experiment selection versus stronger evidence aggregation.

\paragraph{Some models continue querying after sufficient evidence.}
The ledger also measures extra turns after the first strong-evidence turn.
The strongest current models usually stop soon after obtaining sufficient evidence.
On \pgmstruct, Opus-4.6 and Sonnet-4.6 average only about $0.3$ extra turns, Gemini-3-Pro about $0.9$, and Qwen3.6-27B about $0.1$.
By contrast, GPT-OSS-120B averages $3.8$ extra turns, Llama-4-Maverick $6.0$, GPT-4o $2.6$, and Llama-3-70B $8.1$.
This is a separate bottleneck from evidence acquisition: some agents can find useful evidence but fail to recognize that the evidence is already sufficient.

\begin{figure}[t]
    \centering
    \includegraphics[width=0.3\linewidth]{fig4_ledger_bars.pdf}
    \caption{
    Evidence-ledger diagnostics for modular code-using trajectories.
    Left panels report accuracy, evidence acquisition, evidence retention, and ignored evidence.
    Right panels report extra turns after the first strong-evidence turn.
    The dominant variation is in evidence acquisition and stopping behavior; retention is comparatively high once strong evidence appears.
    }
    \label{fig:ledger_bars}
\end{figure}

Together, these results show why trajectory-level evaluation matters.
The same final error can arise from failing to select the right experiment, failing to preserve already-correct evidence, or failing to stop after the answer is supported.
These mechanisms call for different remedies, and they are invisible in final-answer accuracy alone. Full ledger diagnostics are reported in Appendix~\ref{app:ledger_full}.

\section{Limitations and Discussion}
\label{sec:limitations}

Our benchmark is deliberately controlled.
This control is what makes exact scoring, reproducible intervention semantics, and trajectory-level diagnosis possible, but it also limits the kinds of scientific uncertainty represented.
The current worlds use categorical variables and fully observed PGMs, excluding continuous measurements, dose-response curves, temporal dynamics, measurement error, latent confounding, selection effects, missingness, and model misspecification.
The causal-decision dataset is generated from a fixed set of archetypes, so it tests targeted causal decision patterns rather than open-ended scientific discovery.
The results should therefore be interpreted as evidence about active causal investigation in well-specified environments, not as a direct estimate of performance in real laboratories or observational science.

The scenarios are synthetic.
LLMs provide variable names, descriptions, stories, and role plans, while programmatic validation defines the graph, CPDs, interventions, effect-size margins, and gold answers.
This separation prevents LLM-generated text from determining the causal ground truth, but the surface form of the worlds may still reflect the domains, prompts, and model used during generation.
The evaluated instance is also finite: 90 worlds, 240 questions, seven contemporary scientist models, two earlier-generation reference models, and a limited set of agent designs.
The generator can scale to additional worlds, graph sizes, domains, and archetypes without manual answer labeling, but larger public instances are needed to measure robustness across a wider range of scientific settings and model families.

The evidence-ledger analysis is diagnostic rather than definitive.
It relies on an LLM annotator to judge whether a trajectory computed the correct estimand, whether the resulting evidence supports the gold answer, and whether that evidence was retained.
Although the annotator is given the question, gold answer, trajectory context, and code outputs, its judgments may still be imperfect.
The current ledger focuses on modular code-using trajectories, so it should be read as a process-level diagnostic rather than as the primary benchmark score.
We therefore use the ledger to interpret failure mechanisms, while treating final-answer accuracy as the main quantitative metric.
Future versions should expand ledger coverage, calibrate annotations against human judgments, and evaluate richer settings with continuous variables, latent variables, noisy measurements, longer experimental campaigns, and more open-ended scientific goals.

\section{Conclusion}
\label{sec:conclusion}

We introduced \bench, an interactive PGM-based benchmark for evaluating LLM agents as active statistical and causal investigators under exact causal ground truth.
Agents interact with hidden probabilistic graphical worlds through observational and interventional samples, optional code analysis, and budgeted experiment selection, while the evaluator retains oracle-checkable answers for both structural queries and causal decisions.

The main lesson is that scientific-agent evaluation should measure the process of investigation, not only the final answer.
Active evidence collection substantially improves performance over guess-only baselines, showing that models can use interaction to recover hidden statistical and causal structure.
But interaction alone is not enough.
Causal-decision tasks remain more sensitive to whether the agent selects the right constrained causal contrast, and code helps only when the agent can maintain a reliable experiment-analysis loop.

Our evidence-ledger analysis shows that remaining failures are qualitatively different.
Across the evaluated model fleet, evidence retention is usually high once strong evidence appears, while evidence acquisition and stopping behavior remain more diagnostic: some agents never test the right contrast, some ignore earlier strong evidence, and some continue querying after the answer is already supported.
These distinctions matter because they imply different remedies.
Better experiment selection, stronger evidence aggregation, and explicit stopping criteria are separable challenges for AI co-scientists.
\bench provides a controlled and scalable way to evaluate these capabilities under exact causal ground truth, and a step toward scientific agents whose conclusions are supported not only by plausible reasoning, but by evidence they actively gathered and correctly used.

% \begin{ack}
% Use unnumbered first level headings for the acknowledgments. All acknowledgments
% go at the end of the paper before the list of references. Moreover, you are required to declare
% funding (financial activities supporting the submitted work) and competing interests (related financial activities outside the submitted work).
% More information about this disclosure can be found at: \url{https://neurips.cc/Conferences/2026/PaperInformation/FundingDisclosure}.


% Do {\bf not} include this section in the anonymized submission, only in the final paper. You can use the \texttt{ack} environment provided in the style file to automatically hide this section in the anonymized submission.
% \end{ack}

% \section*{References}


\bibliographystyle{unsrtnat}
\bibliography{references}

%%%%%%%%%%%%%%%%%%%%%%%%%%%%%%%%%%%%%%%%%%%%%%%%%%%%%%%%%%%%

\appendix

\section{Additional Benchmark Diagrams}
\label{app:struct_diagram}

\begin{figure}[h]
    \centering
    \includegraphics[width=0.92\linewidth]{diagram-easy.pdf}
    \caption{
    Interaction loop for \pgmstruct.
    The benchmark oracle samples from a hidden PGM, while the LLM scientist agent observes only the scenario, variable metadata, and question.
    The agent requests observational or interventional samples and answers structural statistical or causal queries, such as marginal probabilities, conditional probabilities, independences, or basic intervention effects.
    }
    \label{fig:appendix_struct_overview}
\end{figure}

\section{World Validation Details}
\label{app:validation}

\paragraph{Structural validation.}
For structural worlds, CPD resampling is governed by a total-variation conditional-independence test.
For each tested pair of variables $(A,B)$ and conditioning variable $Z$, we compute
\[
    \mathrm{TV}\!\left(
    P(A,B \mid Z=z),
    P(A \mid Z=z)P(B \mid Z=z)
    \right)
\]
by exact inference for each value $z$ of $Z$.
When the graph implies that $A$ and $B$ should be conditionally dependent given $Z$, we require the distance to exceed a threshold $\varepsilon_{\mathrm{faith}}$.
When the graph implies conditional independence, we require the distance to be below $\varepsilon_{\mathrm{faith}}$.
CPDs are resampled until these checks pass or the attempt budget is exhausted.
Worlds for which no CPD draw passes validation are discarded.
These checks filter accidental independences and unstable dependencies that would make structural questions ambiguous.

\paragraph{Causal-decision validation.}
For causal-decision worlds, each archetype has a dedicated exact-inference validator.
Safety-constrained worlds are rejected unless a safe and effective action exists: at least one intervention improves the target by at least $0.20$ expected-state-index units while changing the safety outcome by at most $0.08$ units.
They are also rejected unless the tempting unsafe action genuinely worsens the safety outcome.

Mediator-structure worlds are rejected unless the total causal effect of treatment on outcome exceeds $0.20$ units and the subvariant-specific signature holds.
For mediated-only worlds, the residual treatment effect after intervening on the mediator must be at most $0.10$ units.
For direct-and-mediated worlds, the residual effect must exceed $0.18$ units.
For not-mediator worlds, fixing the proposed mediator must not substantially attenuate the treatment effect.
For which-mediator worlds, the true mediator's interventional effect must exceed the runner-up by at least $0.20$ units.

Satisficing worlds are rejected unless the improvement distribution has a clean gap that allows the threshold to be placed without ambiguity.
Subgroup-robust worlds are rejected unless the robust intervention improves every subgroup, the average-best intervention fails in at least one subgroup, and the bad intervention harms at least one subgroup.
Invalid-premise worlds are rejected unless the proposed non-intervenable variable has a sufficiently large observational association with the target, and the recommended alternative has a sufficiently large interventional effect.
Across all archetypes, worlds are also rejected when required or forbidden edges cannot be placed within the generated graph, when effect margins are too small, or when the oracle answer is ambiguous.

\section{Structural Question-Type Breakdown}
\label{app:structural_qtype_breakdown}

The structural dataset contains question types that differ in how much experiment selection they require.
Specified marginal, conditional, and independence questions provide the agent with a relatively clear estimand, so active sampling can often estimate the relevant relationship directly.
Set-valued causal-structure questions are harder because the agent must discover and aggregate a set of causal relations across many candidate variables.
This distinction explains why \pgmstruct can show large gains from active sampling while still leaving substantial errors in causal-structure questions.

The fine-grained breakdown should therefore be read as a diagnostic decomposition of \pgmstruct, not merely as a ranking of models.
Independence-style questions mainly test whether the agent can request and interpret the relevant distributional comparison.
Causal-structure questions additionally test whether the agent can plan a set of interventions or comparisons, track multiple candidate variables, and synthesize the resulting evidence into a set-valued answer.
This separation is consistent with the evidence-ledger analysis in Section~\ref{subsec:evidence_ledger_results}: many remaining failures arise from not acquiring the right evidence, rather than from misreading evidence once it is obtained.

\section{Per-category Accuracy on PGM-Struct}
\label{app:structural_category_breakdown}

The structural dataset contains six questions per world that fall into three broad categories: \emph{causal-structure} questions, \emph{marginal-independence} questions, and \emph{conditional-independence} questions.
The aggregate accuracies in Section~\ref{subsec:overall_results} mix these categories.
Figure~\ref{fig:basic_categories} reports accuracy separately for each category, faceted by agent type.

The per-category view exposes a pattern hidden in the dataset-wide average.
Independence questions become much easier once the agent can query the simulator, because the relevant marginal or conditional relationship can be estimated directly from samples.
The remaining error budget on \texttt{PGM-Struct} comes largely from causal-structure questions, especially set-valued enumeration of all causes or all effects of a target variable.
These questions require the agent to discover and aggregate multiple causal relations across candidate variables, rather than estimate one specified relationship.
This separation is consistent with the evidence-ledger results in Section~\ref{subsec:evidence_ledger_results}: the dominant remaining bottleneck is often evidence acquisition, i.e., choosing which causal or statistical contrast to test.

\begin{figure}[t]
    \centering
    \includegraphics[width=\linewidth]{fig6_basic_categories.pdf}
    \caption{
    Per-category accuracy on \texttt{PGM-Struct}, faceted by agent type.
    Bars are grouped by question category; colors and hatches mark the model.
    Causal-structure questions are generally harder than independence questions because they require discovering and aggregating multiple relations rather than estimating a specified relationship.
    }
    \label{fig:basic_categories}
\end{figure}

\section{Overall Accuracy with 95\% Confidence Intervals}
\label{app:accuracy_cis}

Figure~\ref{fig:overall_accuracy_bars} re-presents the overall-accuracy data of Figure~\ref{fig:overall_accuracy} as grouped bars with explicit error bars.
Intervals are $95\%$ percentile bootstrap CIs ($B{=}2000$).
On \texttt{PGM-Struct}, we use a world-clustered bootstrap that resamples the 30 worlds, each containing 6 questions, rather than the 180 individual questions; this preserves the within-world dependence introduced by sharing the same graph and CPDs across questions.
On \texttt{PGM-Decision}, each world contributes one question, so world-level and question-level bootstraps coincide.
For within-world comparisons between agent variants, we use paired bootstrap intervals.

The bar-chart presentation makes two patterns visible.
First, active sampling improves over guess-only answering across models.
Second, code and modular decomposition have model-dependent effects: they help the strongest models, but single-prompt code can underperform active sampling for models that struggle to maintain the long agent loop.
The modular agent often recovers these failures by decomposing planning, analysis, interpretation, and experiment selection into separate stages.

% \begin{figure}[t]
%     \centering
%     \includegraphics[width=0.9\linewidth]{fig1b_combined_accuracy_bars.pdf}
%     \caption{
%     Overall accuracy with $95\%$ world-clustered bootstrap intervals ($B{=}2000$).
%     Same data as Figure~\ref{fig:overall_accuracy}, presented as grouped bars to make per-cell uncertainty visible.
%     Active sampling consistently improves over guess-only answering, while code and modular decomposition have model-dependent effects.
%     }
%     \label{fig:overall_accuracy_bars}
% \end{figure}

\section{Evidence-Ledger Metric Definitions}
\label{app:ledger_definitions}

For trajectory $i$, let $S_i=1$ indicate that the trajectory contains at least one \emph{strong-evidence} turn, and let $C_i=1$ indicate that its final answer is correct.

\paragraph{Evidence acquisition.}
Evidence acquisition is
\[
    \mathrm{Acq}
    =
    \Pr(S_i=1).
\]

\paragraph{Evidence retention.}
Evidence retention is the probability that the final answer is correct after strong evidence has appeared:
\[
    \mathrm{Ret}
    =
    \Pr(C_i=1 \mid S_i=1)
    =
    \frac{\sum_i \mathbf{1}\{S_i=1 \wedge C_i=1\}}
         {\sum_i \mathbf{1}\{S_i=1\}} .
\]

\paragraph{Ignored evidence.}
Ignored evidence is the fraction of failed trajectories that nevertheless contained strong evidence:
\[
    \mathrm{Ignored}
    =
    \Pr(S_i=1 \mid C_i=0)
    =
    \frac{\sum_i \mathbf{1}\{S_i=1 \wedge C_i=0\}}
         {\sum_i \mathbf{1}\{C_i=0\}} .
\]

\paragraph{Extra turns after evidence.}
Extra turns after evidence is the number of turns after the first strong-evidence turn.
It measures whether the agent continues querying after it has already obtained sufficient evidence for the gold answer.
We report its mean and median over trajectories with $S_i=1$.

\section{Full Evidence-Ledger Metrics}
\label{app:ledger_full}

Table~\ref{tab:evidence_ledger_full} reports the full evidence-ledger diagnostics for modular code-using trajectories.
These results are diagnostic rather than primary benchmark scores; the main quantitative metric remains final-answer accuracy.

\begin{table}[h]
\centering
\small
\caption{
Full evidence-ledger diagnostics for modular code-using trajectories.
Acc., Acq., Ret., and Ignored are proportions.
Acq. is evidence acquisition, Ret. is evidence retention, Ignored is the fraction of failed trajectories that previously contained strong evidence, and extra-turn columns report the number of turns after the first strong-evidence turn.
}
\label{tab:evidence_ledger_full}
\resizebox{\textwidth}{!}{
\begin{tabular}{llrrrrrrr}
\toprule
Dataset & Model & $n$ & Acc & Acq. & Ret. & Ignored & Mean extra & Med. extra \\
\midrule
\pgmstruct   & Opus-4.6         & 173 & 0.948 & 0.879 & 0.974 & 0.444 & 0.34 & 0 \\
\pgmstruct   & Sonnet-4.6       & 180 & 0.878 & 0.778 & 0.979 & 0.136 & 0.33 & 0 \\
\pgmstruct   & Gemini-3-Pro     & 180 & 0.939 & 0.872 & 0.968 & 0.455 & 0.87 & 0 \\
\pgmstruct   & DeepSeek-V4-Pro  & 60  & 0.850 & 0.817 & 0.959 & 0.222 & 0.98 & 0 \\
\pgmstruct   & Qwen3.6-27B      & 60  & 0.733 & 0.700 & 0.952 & 0.125 & 0.10 & 0 \\
\pgmstruct   & GPT-OSS-120B     & 60  & 0.800 & 0.780 & 0.920 & 0.330 & 3.80 & 1 \\
\midrule
\pgmstruct   & Llama-4-Maverick & 179 & 0.754 & 0.693 & 0.855 & 0.409 & 5.98 & 1 \\
\pgmstruct   & GPT-4o           & 174 & 0.741 & 0.575 & 0.900 & 0.222 & 2.59 & 1 \\
\pgmstruct   & Llama-3-70B      & 180 & 0.694 & 0.656 & 0.831 & 0.364 & 8.09 & 4 \\
\midrule
\pgmdecision & Opus-4.6         & 40  & 0.925 & 1.000 & 0.925 & 1.000 & 0.42 & 0 \\
\pgmdecision & Sonnet-4.6       & 48  & 0.979 & 0.917 & 1.000 & 0.000 & 0.02 & 0 \\
\pgmdecision & Gemini-3-Pro     & 48  & 0.812 & 0.812 & 0.872 & 0.556 & 2.44 & 1 \\
\pgmdecision & DeepSeek-V4-Pro  & 48  & 0.917 & 0.792 & 0.921 & 0.300 & 0.97 & 0 \\
\pgmdecision & Qwen3.6-27B      & 48  & 0.875 & 0.500 & 1.000 & 0.000 & 0.12 & 0 \\
\pgmdecision & GPT-OSS-120B     & 60  & 0.700 & 0.720 & 0.880 & 0.300 & 1.20 & 0 \\
\midrule
\pgmdecision & Llama-4-Maverick & 48  & 0.604 & 0.333 & 0.875 & 0.105 & 0.88 & 0 \\
\bottomrule
\end{tabular}
}
\end{table}

%%%%%%%%%%%%%%%%%%%%%%%%%%%%%%%%%%%%%%%%%%%%%%%%%%%%%%%%%%%%

\newpage
\input{checklist.tex}

\end{document}